\textbf{C041 抛物线的焦准定义（统一定义e=1）。} $F\notin\ell$，$PF=d(P,\ell)>0$。
$p=d(F,\ell)>0$；顶点是焦点与准线垂足连线的中点。仅在顶点原点、向右开口的标准位置，$F(p/2,0),\ell:x=-p/2$，方程 $y^2=2px$。
非负距离等式一次等价平方即可推导。几何上等价于“经过F并与ℓ相切的圆，其圆心轨迹”。
$PF=x+p/2\ge p/2$，焦半径小于p/2无点，等于p/2只有顶点，大于p/2恰有两个点。
反例：F(3,0)、准线x=5得到顶点(4,0)、左开口 $y^2=-4(x-4)$；不能只看焦点横坐标的正负。
